\begin{scanpage}{0212}{194}
\[E_z^e=\mathcal E_1^{(2)}\delta(z-z_1),\tag{31.03}\]
where
\[\mathcal E_1^{(2)}=\left.\mathcal E^{(2)}(z_1)\right|_{\vartheta=0}
=E\frac{2i\ln\dfrac{i}{\gamma ka}\psi_+^0(L)}{k\sin\vartheta_0\ln\dfrac{2i}{\gamma ka\sin\vartheta_0}}e^{ikL-ikz_2\cos\vartheta_0},\tag{31.04}\]
a wave arises which coincides with wave (30.34) when $\vartheta\approx0$. It is not difficult to see that these external fields actually excite in an infinite single-wire line current waves which are equivalent to the secondary current waves in a passive vibrator [that is, equivalent to those waves which are expressed by the first terms in the brackets of Equation (28.18)]. Therefore, the tertiary waves may be investigated as edge waves radiated by the semi-infinite conductors $(z_1,\infty)$ and $(-\infty,z_2)$ with their excitation by the external fields (31.01) and (31.03), respectively. From Equations (30.25) and (30.32), we find without difficulty the total field radiated with the indicated excitation by the conductor $(z_1,\infty)$
\[E_\vartheta=H_\varphi=\frac{\mathcal E_2^{(2)}}{2\sin\vartheta\ln\dfrac{2i}{\gamma ka\sin\vartheta}}\frac{e^{ikR}}R\bigl[e^{iwz_2}-\psi_-(L)e^{iwz_1+ikL}\bigr]\tag{31.05}\]
and the total field radiated by the conductor $(-\infty,z_2)$
\[E_\vartheta=H_\varphi=\frac{\mathcal E_1^{(2)}}{2\sin\vartheta\ln\dfrac{2i}{\gamma ka\sin\vartheta}}\frac{e^{ikR}}R\bigl[e^{iwz_1}-\psi_+(L)e^{iwz_2+ikL}\bigr].\tag{31.06}\]
As a result, we obtain for the tertiary waves diverging from the ends $z_1$ and $z_2$ the following expressions:
\[E_\vartheta^{(3)}(z_1)=H_\varphi^{(3)}(z_1)=\frac{\mathcal E^{(3)}(z_1)}{2\sin\vartheta\ln\dfrac{2i}{\gamma ka\sin\vartheta}}\frac{e^{ikR}}R e^{-ikz_1\cos\vartheta},\tag{31.07}\]
\[E_\vartheta^{(3)}(z_2)=H_\varphi^{(3)}(z_2)=\frac{\mathcal E^{(3)}(z_2)}{2\sin\vartheta\ln\dfrac{2i}{\gamma ka\sin\vartheta}}\frac{e^{ikR}}R e^{-ikz_2\cos\vartheta},\tag{31.08}\]
where
\end{scanpage}
